\section{Black Hole Solution in Topological Integrity Gravity (TIG)}

\subsection{Metric Ansatz}

We consider a static, spherically symmetric spacetime of the form:

\[
ds^2 = -f(r) dt^2 + f(r)^{-1} dr^2 + r^2 d\Omega^2
\]

In classical General Relativity, the Schwarzschild solution is given by:

\[
f(r) = 1 - \frac{2GM}{r}
\]

\subsection{TIG Modification}

In the TIG framework, we introduce an integrity correction term that modifies the effective gravitational behavior at small scales.

The modified metric function is:

\[
f(r) = 1 - \frac{2GM}{r} \left(1 - e^{-r_0 / r}\right)
\]

where:
\begin{itemize}
\item $G$ is the gravitational constant
\item $M$ is the mass
\item $r_0$ is the fundamental integrity scale
\end{itemize}

\subsection{Properties of the Solution}

\paragraph{Asymptotic Behavior}

For large $r$, the exponential term vanishes:

\[
e^{-r_0 / r} \rightarrow 1
\]

and we recover the Schwarzschild solution:

\[
f(r) \approx 1 - \frac{2GM}{r}
\]

\paragraph{Regular Core}

For $r \rightarrow 0$, the correction term suppresses the singularity:

\[
f(r) \rightarrow 1
\]

This indicates the absence of a classical curvature singularity.

\subsection{Physical Implications}

The TIG modification leads to:

\begin{itemize}
\item Regularization of the black hole core
\item Modified horizon structure
\item Potential deviations in black hole thermodynamics
\end{itemize}

In particular, the event horizon is defined by:

\[
f(r_H) = 0
\]

which must be solved numerically in the TIG framework.

\subsection{Interpretation}

The exponential correction encodes the effect of topological integrity, acting as a stabilizing mechanism that prevents the formation of singularities.

This suggests that gravitational collapse is regulated by an intrinsic geometric principle rather than leading to divergence.
